% from ega4-8-fr.tex
\oldpage[IV-3]{15}
Inversement, pour montrer que $u_\lambda$ est ouvert lorsque tous les
$u_{\lambda\mu}$ le sont pour $\mu\geq\lambda$, il suffit de voir que pour
tout ouvert quasi-compact $V$ de $S$, $u_\lambda(V)$ est ouvert dans
$S_\lambda$~; mais il existe alors $\mu\geq\lambda$ tel que
$V=u_\mu^{-1}(V_\mu)$, où $V_\mu$ est ouvert dans $S_\mu$ (8.2.11)~; comme
$u_\mu$ est surjectif, on a $V_\mu=u_\mu(V)$ et
$u_\lambda(V)=u_{\lambda\mu}(u_\mu(V))$ est donc ouvert par hypothèse.

On notera qu'il peut se faire que tous les $u_{\lambda\mu}$ soient ouverts
sans qu'il en soit ainsi de $u_\lambda$ si les $u_\mu$ ne sont pas
surjectifs. On en a un exemple en considérant un anneau intègre $A$ qui n'est
pas un corps, et son corps des fractions $K$, qui est limite inductive des
$A_f$, où $f$ parcourt $A-\{0\}$~; si l'on pose
$S_1=\operatorname{Spec}(A)$, $S_f=\operatorname{Spec}(A_f)$, on a
$S=\varprojlim S_f=\operatorname{Spec}(K)$, et le morphisme $S\to S_1$ n'est
pas ouvert, bien que les morphismes $S_f\to S_1$ le soient.

\begin{env}[8.3.9]
\label{IV.8.3.9-fr}
Pour tout préschéma $X$, nous noterons comme d'ordinaire $\mathfrak P(X)$
l'ensemble des parties de l'ensemble sous-jacent à $X$, par $\mathfrak C(X)$
(resp. $\mathfrak D_c(X)$, $\mathfrak F_c(X)$, $\mathfrak{LF}_c(X)$)
l'ensemble des parties constructibles (resp. constructibles et ouvertes,
resp. constructibles et fermées, resp. constructibles et localement fermées)
de $X$. Il est clair que
$(\mathfrak P(S_\lambda),u_{\lambda\mu}^{-1})$ est un \emph{système inductif
d'ensembles} et que les applications
$u_\lambda^{-1}:\mathfrak P(S_\lambda)\to\mathfrak P(S)$ forment un système
inductif d'applications, d'où, en passant à la limite inductive, une
application canonique
\begin{equation}
u_{\mathfrak P}:\varinjlim\mathfrak P(S_\lambda)\to\mathfrak P(S).
\tag{8.3.9.1}\label{IV.8.3.9.1-fr}
\end{equation}
En outre, il résulte de (1.8.2) que $u_{\lambda\mu}^{-1}$ applique
$\mathfrak C(S_\lambda)$ (resp. $\mathfrak D_c(S_\lambda)$,
$\mathfrak F_c(S_\lambda)$, $\mathfrak{LF}_c(S_\lambda)$) dans
$\mathfrak C(S_\mu)$ (resp. $\mathfrak D_c(S_\mu)$,
$\mathfrak F_c(S_\mu)$, $\mathfrak{LF}_c(S_\mu)$) et que
$u_\lambda^{-1}$ applique $\mathfrak C(S_\lambda)$ (resp.
$\mathfrak D_c(S_\lambda)$, $\mathfrak F_c(S_\lambda)$,
$\mathfrak{LF}_c(S_\lambda)$) dans $\mathfrak C(S)$ (resp.
$\mathfrak D_c(S)$, $\mathfrak F_c(S)$, $\mathfrak{LF}_c(S)$). On a donc par
restriction de (8.3.9.1), des applications canoniques
\begin{align}
\varinjlim\mathfrak C(S_\lambda)&\longrightarrow\mathfrak C(S)
\tag{8.3.9.2}\label{IV.8.3.9.2-fr}\\
\varinjlim\mathfrak D_c(S_\lambda)&\longrightarrow\mathfrak D_c(S)
\tag{8.3.9.3}\label{IV.8.3.9.3-fr}\\
\varinjlim\mathfrak F_c(S_\lambda)&\longrightarrow\mathfrak F_c(S)
\tag{8.3.9.4}\label{IV.8.3.9.4-fr}\\
\varinjlim\mathfrak{LF}_c(S_\lambda)&\longrightarrow\mathfrak{LF}_c(S).
\tag{8.3.9.5}\label{IV.8.3.9.5-fr}
\end{align}
\end{env}

\begin{env}[8.3.10]
\label{IV.8.3.10-fr}
Soit $g_\alpha:X_\alpha\to S_\alpha$ un morphisme~; avec les notations de
(8.2.5) on a comme ci-dessus une application canonique
$v_{\mathfrak P}:\varinjlim\mathfrak P(X_\lambda)\to\mathfrak P(X)$~;
d'autre part, on a des morphismes de projection
$g_\lambda:X_\lambda\to S_\lambda$ pour tout $\lambda\geq\alpha$ et un
morphisme de projection $g:X\to S$. Il est clair que les
$g_\lambda^{-1}:\mathfrak P(S_\lambda)\to\mathfrak P(X_\lambda)$ forment un
système inductif d'applications, et que les diagrammes
\[
\xymatrix{
  \mathfrak P(S_\lambda) \ar[r]^{u_\lambda^{-1}} \ar[d]_{g_\lambda^{-1}}
    & \mathfrak P(S) \ar[d]^{g^{-1}} \\
  \mathfrak P(X_\lambda) \ar[r]_{v_\lambda^{-1}} & \mathfrak P(X)
}
\]
